\subsection{Historical coverage models}
\label{app:method-evidence}

\paragraph{Pair coverage and evidence selection.}
For each target--prior pair $(x,p_k)$, the pair coverage model estimates how
much of the target the prior accounts for. It outputs a score $s_k\in[0,1]$
and a distribution over four classes: unrelated, related but uncovered,
partially covered, and largely covered. A cross-encoder predicts relatedness
$a_k$, coverage conditional on relatedness $b_k$, and large coverage
conditional on coverage $d_k$. These probabilities define the class
distribution
\begin{equation}
 \boldsymbol{\pi}^{\mathrm{pair}}_k
 = [1-a_k,\;a_k(1-b_k),\;a_kb_k(1-d_k),\;a_kb_kd_k].
 \label{eq:pair-probs}
\end{equation}
With classes indexed from $0$ to $3$, the coverage score is
$s_k=\tfrac12\pi^{\mathrm{pair}}_{k,2}+\pi^{\mathrm{pair}}_{k,3}$.
The reference pipeline ranks temporally admissible candidates by
$s_k+0.02\ell_k/3$, where $\ell_k$ is the predicted ordinal coverage label,
and retains at most $K$ priors. The reported IdeaDelta checkpoint instead
uses the ordering and pair annotations supplied in the dataset ranking
records, as detailed under the evidence-source protocol below. In both
cases, the retained priors and their pair coverage features form the evidence
inputs to the subsequent stages.

\paragraph{Joint coverage.}
Different priors may cover different parts of the target. We therefore
estimate their joint coverage with a multiple-instance model, treating each
pair $(x,p_k)$ as an instance and the retained evidence set as a bag
\citep{ilse2018attentionmil}. Let $n=|\mathcal{E}_x|$ denote the number of
retained priors. A shared cross-encoder maps each pair to a representation
$\boldsymbol{h}^{J}_k$, from which the model computes attention weights
$\beta_k$ and internal evidence scores $e_k$:
\begin{equation}
 \beta_k=\operatorname{softmax}_{k}(\boldsymbol{w}_a^\top\boldsymbol{h}^{J}_k+b_a),
 \qquad e_k=\sigma(\boldsymbol{w}_e^\top\boldsymbol{h}^{J}_k+b_e).
\end{equation}
For a nonempty evidence set, noisy-OR pools the evidence scores as
\begin{equation}
 e_{\cup}=1-\prod_{k=1}^{n}(1-e_k)
 \label{eq:noisy-or}
\end{equation}
The model summarizes these scores by their noisy-OR aggregate, maximum,
mean, and the fraction of available evidence slots. It concatenates these
statistics with attention-pooled and elementwise maximum-pooled pair
representations:
\begin{align}
 \boldsymbol{v}_E
 &=\left[e_{\cup},\;
          \max_k e_k,\;\frac{1}{n}\sum_{k=1}^{n}e_k,\;\frac{n}{K}\right],
 \\
 \boldsymbol{z}_J
 &=f_J\left(\left[\sum_{k=1}^{n}\beta_k\boldsymbol{h}^{J}_k;
              \operatorname{max}_{k}\boldsymbol{h}^{J}_k;
              \boldsymbol{v}_E\right]\right).
 \label{eq:joint-bag}
\end{align}
Here $f_J$ is an MLP. Attention pooling aggregates
evidence across priors, while maximum pooling retains the strongest feature
values. A hierarchical head maps $\boldsymbol{z}_J$ to three probabilities:
any coverage, substantial coverage conditional on any coverage, and mostly
covered conditional on substantial coverage. Using the factorization in
Eq.~(\ref{eq:pair-probs}), these probabilities yield
$\boldsymbol{\pi}^{J}$ over four joint-coverage classes: \emph{not},
\emph{weakly}, \emph{partially}, and \emph{mostly} covered. The scalar
joint-coverage estimate is
$C_J=\sum_{\ell=0}^{3}\mu_\ell\pi^J_\ell$, where $\mu_\ell$ is the mean
annotated coverage score for class $\ell$ on the training split.
Both $C_J$ and $\boldsymbol{\pi}^{J}$ are supplied as fixed inputs to the
residual and novelty modules. The noisy-OR aggregate $e_{\cup}$ contributes
one of the four statistics in $\boldsymbol{v}_E$; the joint-coverage
prediction is learned from the full bag representation.

\subsection{Residual representations and pooling}
\label{app:method-representations}

\paragraph{Pair encoding.}
For each target--prior pair $(x,p_k)$, the residual cross-encoder produces
token states $\boldsymbol{H}_k\in\mathbb{R}^{L\times d}$, where $L$ is the
input sequence length. We apply masked mean pooling to obtain the pair
representation and augment it with projected pair coverage features:
\begin{equation}
 \boldsymbol{v}_k=\operatorname{MeanPool}(\boldsymbol{H}_k),
 \qquad
 \widetilde{\boldsymbol{v}}_k=\boldsymbol{v}_k+f_P(\boldsymbol{f}_k),
\end{equation}
where $f_P$ consists of a linear map, layer normalization, GELU, and dropout.
The residual encoder shares parameters across all retained priors. To
construct the target token bank, we average target-token states at matching
positions across valid pairs. Masked mean pooling over this bank yields the
target representation $\hbar_x$ used in residual fusion.

\paragraph{Coverage reconstruction.}
To supervise the unit coverage estimates, we aggregate them into a scalar
joint-coverage prediction. The pooling weights depend on both the semantic
representation and the activity of each unit:
\begin{equation}
 \widehat C_U=\sum_{m=1}^{M}\omega_m c_m,\qquad
 \omega_m=\frac{v_m\exp(\boldsymbol{w}_c^\top\boldsymbol{u}_m+b_c)}
 {\sum_{r=1}^{M}v_r\exp(\boldsymbol{w}_c^\top\boldsymbol{u}_r+b_c)}.
 \label{eq:reconstructed-coverage}
\end{equation}
The resulting estimate $\widehat C_U$ is supervised by annotated coverage
and the fixed Stage~A prediction. This aggregate supervision constrains unit
coverage but does not uniquely determine the latent support matrix. We use
noisy-OR as a differentiable rule for aggregating support; its use does not
establish that support from different priors is statistically independent.

\paragraph{Interaction encoding and attention.}
For each unit pair $i<j$, an MLP $f_I$ encodes the two units and their
elementwise comparisons into an interaction representation. A sigmoid head
then predicts the interaction strength:
\begin{equation}
 \boldsymbol{t}_{ij}=f_I([\boldsymbol{u}_i;\boldsymbol{u}_j;
 \boldsymbol{u}_i\odot\boldsymbol{u}_j;|\boldsymbol{u}_i-\boldsymbol{u}_j|]),
 \qquad g_{ij}=\sigma(\boldsymbol{w}_g^\top\boldsymbol{t}_{ij}+b_g).
\end{equation}
Using the interaction residual weights $r_{ij}$ defined in Stage~B, we
compute attention over all $\binom{M}{2}$ unit pairs:
\begin{equation}
 \eta_{ij}=\operatorname{softmax}_{i<j}
 \left(\boldsymbol{w}_I^\top\boldsymbol{t}_{ij}+b_I+
 \log(r_{ij}+\varepsilon)\right),\qquad
 \boldsymbol{z}_I=\sum_{i<j}\eta_{ij}r_{ij}\boldsymbol{t}_{ij}.
 \label{eq:interaction-pooling}
\end{equation}
The logarithmic bias directs attention toward pairs with larger residual
weights. Multiplication by $r_{ij}$ during pooling also preserves the scale
of the residual weights, allowing weak residuals to contribute less even
after attention normalization.

\subsection{Residual fusion, prediction heads, and decision rules}
\label{app:method-heads}

\paragraph{Context gate.}
The fusion gate conditions the contribution of global context on the two
residual representations and the uncertainty of the joint-coverage
prediction. Let $H_J$ denote the normalized entropy of
$\boldsymbol{\pi}^{J}$. The gate is
\begin{equation}
 \gamma=\sigma\left(f_\gamma([
 \boldsymbol{z}_U;\boldsymbol{z}_I;\boldsymbol{z}_G;H_J])\right).
\end{equation}
The gated context $\gamma\boldsymbol{z}_G$ enters residual fusion alongside
$\boldsymbol{z}_U$ and $\boldsymbol{z}_I$. Global context
$\boldsymbol{z}_G$ also enters the final shared representation directly.
Thus, $\gamma$ controls the context contribution within the residual branch,
while the direct context path remains available to the prediction heads.

\paragraph{Auxiliary predictions.}
The auxiliary heads use the shared representation $\boldsymbol{h}$ to predict
substantive difference $\widehat d$, surface change $\widehat s$, evidence
sufficiency $\widehat e$, and residual-type probabilities
$\boldsymbol{\pi}^{T}$. The first three outputs are sigmoid scores; the type
prediction uses the hierarchy described below.

\paragraph{Residual types.}
The type hierarchy distinguishes a minor variant from six substantive families:
conceptual innovation, mechanism or system, task or application, training or
evaluation, data-related innovation, and combination. Conditional fine heads
yield eleven types in total. The minor/substantive gate additionally reads
$\widehat d$, $\widehat s$, and $\max_{i<j}r_{ij}$. The probability assigned
to the combination type is multiplied by
$0.70+0.30\sigma(f_{\mathrm{comb}}(\boldsymbol{z}_I))$ before the hierarchical
distribution is renormalized. The final type distribution is a convex mixture
of the hierarchical and flat predictions, with flat weight $0.15$.

\paragraph{Novelty branches.}
All novelty branches receive the prediction features
\begin{equation}
 \boldsymbol{\xi}=[\boldsymbol{h};\widehat d;\widehat s;\widehat e;
                    \boldsymbol{\pi}^{T};C_J;\boldsymbol{\pi}^{J}].
\end{equation}
We encode \textsc{low}, \textsc{weak}, \textsc{moderate}, and \textsc{high}
as $0$, $1$, $2$, and $3$, respectively. A direct linear classifier and a
two-layer MLP classifier map $\boldsymbol{\xi}$ to class distributions
$\boldsymbol{p}^{D}$ and $\boldsymbol{p}^{R}$. The MLP branch uses the
full prediction features $\boldsymbol{\xi}$.

The ordinal branch models cumulative probabilities of reaching each grade
threshold. It produces three sigmoid gates $g_1,g_2,g_3$ from
$\boldsymbol{\xi}$ and defines
\begin{equation}
 q_1=g_1,\quad q_2=g_1g_2,\quad q_3=g_1g_2g_3,
 \qquad \boldsymbol{p}^{O}=[1-q_1,\;q_1-q_2,\;q_2-q_3,\;q_3].
\end{equation}
Here $q_\ell$ estimates the probability that $y\geq\ell$, with
$q_1\geq q_2\geq q_3$ by construction.

\paragraph{Probability fusion.}
We aggregate the three branch distributions through a convex mixture:
\begin{equation}
 \boldsymbol{p}=\lambda_D\boldsymbol{p}^{D}
               +\lambda_O\boldsymbol{p}^{O}
               +\lambda_R\boldsymbol{p}^{R},\qquad
 \lambda_D+\lambda_O+\lambda_R=1,
 \quad \lambda_D,\lambda_O,\lambda_R\geq0.
 \label{eq:novelty-fusion}
\end{equation}
The fused distribution provides the initial grade for the refinement below.

\paragraph{Moderate-to-high refinement.}
A separate binary head $h_{MH}=\sigma(f_{MH}(\boldsymbol{\xi}))$ scores
the distinction between moderate and high novelty. At inference, let
$y_0=\arg\max_\ell p_\ell$ for the mixture in
Eq.~(\ref{eq:novelty-fusion}), and define
\begin{equation}
 h_{\mathrm{high}}=\alpha_Dp^D_3+\alpha_Op^O_3+\alpha_{MH}h_{MH},\qquad
 \widehat y=
 \begin{cases}
 3,&y_0=2\ \text{and}\ h_{\mathrm{high}}\geq\theta_{MH},\\
 y_0,&\text{otherwise}.
 \end{cases}
\end{equation}
where $\alpha_D,\alpha_O,\alpha_{MH}\geq0$ and
$\alpha_D+\alpha_O+\alpha_{MH}=1$. The refinement applies a threshold to
this weighted score and does not require separate agreement from all three
heads. For example, an initial \textsc{moderate} prediction is promoted to
\textsc{high} if the combined score reaches $\theta_{MH}$, even if the
individual heads favor different grades. The distribution-mixture weights
and refinement threshold are selected on development data and fixed for
testing. Appendix~\ref{app:method-training} reports the implementation values.

\paragraph{Continuous outputs.}
A separate sigmoid head predicts a continuous novelty score. At inference,
novelty grades are determined by probability fusion and the moderate-to-high
refinement; the continuous score output does not enter that decision rule.
The continuous score head and the substantive-difference head are separate
readouts; the latter is used in the blockwise RINoBench evaluation.

\subsection{Training and implementation details}
\label{app:method-training}

This section supplements the training objective in
Eq.~(\ref{eq:training-objective}) with the loss definitions and implementation
settings for the prediction modules of Appendix~\ref{app:method-heads} and the
residual decomposition of Appendix~\ref{app:method-representations}.

\paragraph{Supervision.}
The classification group $\mathcal{L}_{\mathrm{nov}}$ trains the direct and MLP
classifiers with class-weighted cross-entropy, the ordinal outputs against
$\mathbf{1}[y\geq\ell]$, $\ell=1,2,3$, with weighted binary cross-entropy, and
the fused distribution with a negative log-likelihood. The moderate-to-high head
sees only moderate and high examples and adds a positive--negative margin to its
binary cross-entropy, weighting moderate examples scored as high more heavily.
The type group $\mathcal{L}_{\mathrm{type}}$ combines the hierarchy's
conditional losses, a flat classification loss, binary supervision for
nontrivial combinations, and consistency penalties tying the combination branch
to the type predictions. The separate novelty score regresses the continuous
annotation target
\begin{equation}
 s^*=(1-C^*)d^*e^*(1-0.45s^*_{\mathrm{surface}}),
\end{equation}
formed from annotated coverage, substantive difference, evidence sufficiency,
and surface change.

\paragraph{Scalar heads: level and ranking terms.}
Let $\ell_\delta$ denote smooth-$L_1$ loss and $w_x$ the example weight. For
$a\in\mathcal{A}$, the level loss is
$\mathcal{L}_a=\sum_xw_x\ell_\delta(\widehat a_x-a^*_x)/\sum_xw_x$, where
$a^*_x$ is the annotation and $\widehat a_x$ the prediction. The corresponding
ranking term takes within-minibatch pairs separated by
$|a^*_x-a^*_{x'}|\geq0.05$ and contributes
$[0.02-\operatorname{sgn}(a^*_x-a^*_{x'})(\widehat a_x-\widehat a_{x'})]_+$,
weighted by $\sqrt{w_xw_{x'}}$; it is dropped when a minibatch holds no
eligible pair. The ranking terms of $\Delta$, $S$, $E$, and $C$ are summed into
$\mathcal{L}_{\mathrm{rank}}$ with the internal weights of
Table~\ref{tab:method-settings}; the novelty score's ranking term is carried
inside $\mathcal{L}_{\nu}$ instead, where it takes most of the weight because
the grade is read off the score by thresholding and therefore depends only on
order.

The ranking terms are not redundant with the level losses. At the transition
$\delta$ and target spread used here the level losses operate in their $L_1$
regime, whose optimum is the median of the target, and a head supervised by the
level loss alone drifts towards emitting a near-constant value; the evidence
head is the clearest case. Ranking is scale-free and counteracts this.

\paragraph{Coverage term.}
Beyond its level loss against the annotated joint value $C^*$,
$\mathcal{L}_{C}$ adds a consistency loss against the upstream score $C_J$ at
relative weight $0.2$, reweighted by $1-H_J$ so that examples with uncertain
upstream coverage count for less. No observed unit-level support matrix is
supplied as a target.

\paragraph{Decomposition regularization and variance floor.}
$\mathcal{L}_{\mathrm{dec}}$ combines three terms. Attention regularization
penalizes the squared entries of $\widetilde\alpha\widetilde\alpha^\top-I$ for
the row-normalized attention matrix $\widetilde\alpha$, plus half the mean
off-diagonal cosine similarity. Query diversity penalizes the mean squared
cosine similarity between distinct queries. A usage penalty
$D_{\mathrm{KL}}(\boldsymbol{\rho}\|\boldsymbol{u})$, with
$\rho_m=\sum_jB_{mj}/\sum_{r,j}B_{rj}$ and $\boldsymbol{u}$ uniform over the $M$
units, keeps units from going unused. Their relative weights are given in
Table~\ref{tab:method-settings}.

The variance floor is a separate term rather than a component of
$\mathcal{L}_{\mathrm{dec}}$, and enters
Eq.~(\ref{eq:training-objective}) under its own coefficient:
\begin{equation}
 \mathcal{L}_{\mathrm{var}}
 =\mathbb{E}_x\bigl[[\nu-\operatorname{Var}_m(c_{x,m})]_+/\nu\bigr].
\end{equation}
It acts on the unit coverages rather than on the token assignments, and keeps
the $c_m$ within an example from coinciding. The reported configuration uses
softmax competition with dot-product scores.

\paragraph{Curriculum and feature perturbation.}
For a randomly sampled fraction of training examples the joint-coverage features,
together with the associated probability and entropy features, are replaced by
values derived from annotated joint coverage, with probability $0.50$ in epoch
1, $0.20$ in epoch 2, and $0$ thereafter. This is training-time supervision
only: development and test runs always use the upstream joint predictions.
Selected joint-coverage features on the global-context path also receive
Gaussian noise of standard deviation $0.03$ and independent feature dropout at
probability $0.10$, which leaves untouched both the direct coverage inputs to
the novelty heads and the entropy input to the fusion gate.

\paragraph{Inputs and initialization.}
Pair probabilities missing from a supplied ranking record are reconstructed by
the input adapter from the available coverage label and score; pair labels and
scores also appear in the pair text input. The two facet embeddings entering
$\boldsymbol{z}_G$ have dimension $32$ each. The model uses DeBERTa-v3-base
with hidden dimension $768$, maximum pair-input length $384$, $K=5$, and $M=8$;
the residual encoder is initialized from the learned joint model and upstream
checkpoints stay fixed. Padded priors have zero support, support values are
capped at $0.995$ for numerical stability, and an empty evidence set falls back
to a target-containing input with all prior support masked out. Fusion uses
refinement weights $(\alpha_D,\alpha_O,\alpha_{MH})=(0.40,0.30,0.30)$ and
mixture weights $(\lambda_D,\lambda_O,\lambda_R)=(0.50,0.35,0.15)$ in training;
development-set calibration selects $(0.60,0.20,0.20)$ and $\theta_{MH}=0.32$
for inference, fixed thereafter. Table~\ref{tab:method-settings} lists the
remaining settings.

\begin{table}[t]
\centering
\small
\caption{Residual-stage settings for the reported IdeaDelta checkpoint.
The checkpoint was selected on development data at epoch 5 of 8.}
\label{tab:method-settings}
\begin{tabular}{@{}lr@{}}
\toprule
Setting & Value \\
\midrule
Training batch size / evaluation batch size & 2 / 8 \\
Head learning rate / encoder learning rate & $1.5\times10^{-5}$ / $3\times10^{-6}$ \\
Weight decay / dropout & 0.01 / 0.15 \\
Smooth-$L_1$ transition $\delta$ & 0.25 \\
\addlinespace
$\lambda_\Delta$ / $\lambda_S$ / $\lambda_E$ & 1.0 / 0.6 / 0.8 \\
$\lambda_C$ / $\lambda_\nu$ & 0.3 / 0.6 \\
$\lambda_{\mathrm{nov}}$ / $\lambda_{\mathrm{type}}$ & 0.8 / 1.0 \\
$\lambda_{\mathrm{rank}}$ & 0.5 \\
$\lambda_{\mathrm{dec}}$ / $\lambda_{\mathrm{var}}$ & 0.05 / 0.02 \\
\addlinespace
Within $\mathcal{L}_{\mathrm{rank}}$: $\Delta$ / $S$ / $E$ / $C$ & 1.0 / 0.5 / 0.5 / 0.5 \\
Within $\mathcal{L}_{\nu}$: ranking weight & 2.0 \\
Within $\mathcal{L}_{C}$: Stage~A consistency weight & 0.2 \\
Within $\mathcal{L}_{\mathrm{dec}}$: attention / query / usage & 0.7 / 0.1 / 0.2 \\
Variance floor $\nu$ & 0.0025 \\
\addlinespace
Seed & 173 \\
\bottomrule
\end{tabular}
\end{table}

\paragraph{Evidence-source protocol.}
The reference pair pipeline predicts coverage features and ranks priors by the
rule in Appendix~\ref{app:method-evidence}, while the upstream joint model is
trained with bag-level coverage supervision and an auxiliary pair-score
regression objective. For the reported checkpoint the evidence order and pair
annotations come from the dataset ranking records, so the configuration
evaluates novelty prediction with supplied pair annotations and predicted joint
coverage. Applying the model to an unseen idea additionally requires computing
the pair evidence features.
